\documentclass[12pt, a4paper]{article}
\usepackage[utf8]{inputenc}
\usepackage[T1]{fontenc}
\usepackage[ngerman]{babel}
\usepackage{amsmath, amssymb, amsthm}
\usepackage{geometry}
\usepackage{graphicx}
\usepackage{xcolor}
\usepackage{tcolorbox}
\usepackage{enumitem}
\usepackage{listings}
\usepackage{fancyhdr}
\usepackage{titlesec}
\usepackage{mdframed}
\usepackage{tikz}
\usepackage{pgfplots}
\pgfplotsset{compat=1.18}

\geometry{a4paper, margin=2.5cm}

\definecolor{theoryblue}{RGB}{30, 90, 160}
\definecolor{exerciseorange}{RGB}{200, 80, 0}
\definecolor{solutiongreen}{RGB}{0, 120, 60}
\definecolor{hintgray}{RGB}{100, 100, 100}
\definecolor{codebg}{RGB}{245, 245, 245}

% Theory box
\tcbuselibrary{skins, breakable}
\newtcolorbox{theorybox}[1]{
  colback=theoryblue!5,
  colframe=theoryblue,
  fonttitle=\bfseries\color{white},
  title=#1,
  breakable
}

% Exercise box
\newtcolorbox{exercisebox}[1]{
  colback=exerciseorange!5,
  colframe=exerciseorange,
  fonttitle=\bfseries\color{white},
  title=Aufgabe #1,
  breakable
}

% Solution space box
\newtcolorbox{solutionspace}[1][6cm]{
  colback=white,
  colframe=black!20,
  fonttitle=\itshape\color{hintgray},
  title=Rechenraum,
  height=#1,
  breakable
}

% Hint box
\newtcolorbox{hintbox}{
  colback=solutiongreen!5,
  colframe=solutiongreen,
  fonttitle=\bfseries\color{white},
  title=Hinweis,
  breakable
}

% Julia code style
\lstdefinelanguage{Julia}{
  keywords={function, end, return, if, else, elseif, for, while, in, begin, let, do, try, catch, finally, using, import, export, mutable, struct, abstract, type},
  keywordstyle=\color{theoryblue}\bfseries,
  comment=[l]{\#},
  commentstyle=\color{hintgray}\itshape,
  stringstyle=\color{solutiongreen},
  morestring=[b]",
  sensitive=true
}
\lstset{
  language=Julia,
  backgroundcolor=\color{codebg},
  basicstyle=\ttfamily\small,
  breaklines=true,
  frame=single,
  rulecolor=\color{black!30},
  numbers=left,
  numberstyle=\tiny\color{hintgray}
}

\pagestyle{fancy}
\fancyhf{}
\rhead{\textcolor{hintgray}{\small Rocket Science Workbook}}
\lhead{\textcolor{hintgray}{\small \leftmark}}
\cfoot{\thepage}

\titleformat{\section}{\large\bfseries\color{theoryblue}}{}{0em}{}[\titlerule]
\titleformat{\subsection}{\normalsize\bfseries\color{exerciseorange}}{}{0em}{}


\begin{document}

\begin{center}
  {\LARGE \textbf{Chapter 03: The ISA Atmosphere}}\\[0.5em]
  {\large \textcolor{hintgray}{Deriving Temperature, Pressure and Density from First Principles}}\\[0.3em]
  {\small \textcolor{hintgray}{Rocket Science Workbook — simulation-2D-jl}}
\end{center}
\vspace{1em}\hrule\vspace{1.5em}

% ──────────────────────────────────────────────────────────────
\section{Why the Atmosphere Matters}

\begin{theorybox}{Drag Force}
\[
  F_\text{drag} = \frac{1}{2}\,\rho(h)\,v^2\,C_d\,A
\]
The air density $\rho(h)$ drops by roughly 3 orders of magnitude from sea level to 80~km.
A realistic $\rho(h)$ model is essential — the simple exponential used in the 1D simulation is only valid below ~15~km.
\end{theorybox}

% ──────────────────────────────────────────────────────────────
\section{The Hydrostatic Equation: Full Derivation}

Consider a thin horizontal layer of atmosphere at altitude $h$ with thickness $dh$:

\begin{itemize}
  \item Volume of layer per unit area: $dV = 1 \cdot dh$
  \item Mass of layer: $dm = \rho(h)\,dh$
  \item Weight of layer (force on air below): $dF = \rho(h)\,g\,dh$
  \item This weight must be supported by a pressure difference: $dp = -dF/A = -\rho g\,dh$
\end{itemize}

\begin{theorybox}{Hydrostatic Equation (ODE for pressure)}
\[
  \boxed{\frac{dp}{dh} = -\rho(h)\,g}
\]
This is a first-order ODE. To solve it, we need $\rho$ as a function of $p$ and $T$ — provided by the ideal gas law.
\end{theorybox}

\subsection*{Ideal Gas Law}

\[
  p = \rho\,R_s\,T \quad \Rightarrow \quad \rho = \frac{p}{R_s\,T}
\]

where $R_s = 287.058\;\text{J/(kg\,K)}$ is the specific gas constant for dry air.
Substituting into the hydrostatic equation:

\[
  \frac{dp}{dh} = -\frac{p}{R_s\,T(h)}\,g
\]

% ──────────────────────────────────────────────────────────────
\section{ISA Layers and Analytical Solutions}

\subsection*{Layer 1: Troposphere ($0 \leq h \leq 11\,000\;\text{m}$)}

\begin{theorybox}{Temperature Profile (given empirically)}
\[
  T(h) = T_0 - L\,h, \quad T_0 = 288.15\;\text{K},\quad L = 0.0065\;\text{K/m}
\]
\end{theorybox}

Substituting $T(h) = T_0 - Lh$ into the hydrostatic ODE:
\[
  \frac{dp}{p} = -\frac{g}{R_s(T_0 - Lh)}\,dh
\]

Integrate both sides (left: $\int dp/p = \ln p$; right: requires substitution $u = T_0 - Lh$):

\[
  \ln\frac{p}{p_0} = -\frac{g}{R_s L}\ln\frac{T_0-Lh}{T_0} = \frac{g}{R_s L}\ln\frac{T_0}{T_0-Lh}
\]

\begin{theorybox}{Troposphere: Analytical Solution}
\[
  p(h) = p_0 \left(\frac{T_0 - Lh}{T_0}\right)^{\!\frac{g}{R_s L}}
         = p_0 \left(\frac{T(h)}{T_0}\right)^{5.2561}
\]
\[
  \rho(h) = \rho_0 \left(\frac{T(h)}{T_0}\right)^{\!\frac{g}{R_s L}-1}
           = \rho_0 \left(\frac{T(h)}{T_0}\right)^{4.2561}
\]
The exponent $\dfrac{g}{R_s L} = \dfrac{9.80665}{287.058 \times 0.0065} \approx 5.2561$.
\end{theorybox}

\subsection*{Layer 2: Lower Stratosphere ($11\,000 \leq h \leq 20\,000\;\text{m}$)}

\begin{theorybox}{Isothermal Layer: $T = T_{11} = 216.65\;\text{K}$ (constant)}
The hydrostatic ODE simplifies to:
\[
  \frac{dp}{dh} = -\frac{g}{R_s\,T_{11}}\,p \quad \Rightarrow \quad \frac{dp}{p} = -\lambda\,dh,
  \quad \lambda = \frac{g}{R_s\,T_{11}}
\]
This is exactly $\dot{y} = -\lambda y$ — the decay ODE from Chapter 01!\\
Solution by separation:
\[
  p(h) = p_{11}\,\exp\!\left(-\frac{g}{R_s\,T_{11}}(h - 11000)\right)
\]
\end{theorybox}

% ──────────────────────────────────────────────────────────────
\section{Speed of Sound and Mach Number}

\begin{theorybox}{Acoustic Wave Speed in an Ideal Gas}
\[
  c(h) = \sqrt{\gamma\,R_s\,T(h)}, \qquad \gamma = 1.4
\]
At sea level: $c_0 = \sqrt{1.4 \times 287.058 \times 288.15} \approx 340.3\;\text{m/s}$

At the tropopause ($T_{11} = 216.65\;\text{K}$): $c_{11} = \sqrt{1.4 \times 287 \times 216.65} \approx 295.1\;\text{m/s}$

The Mach number: $Ma = v / c(h)$ — same speed gives \textit{higher} Mach at altitude!
\end{theorybox}

% ──────────────────────────────────────────────────────────────
\section{Mach-Dependent Drag Coefficient}

\begin{theorybox}{Transonic Drag Rise}
Near $Ma=1$, shock waves form on the rocket body, dramatically increasing drag.
We model this with a Gaussian peak:
\[
  C_d(Ma) = C_{d,\infty} + \Delta C_d \cdot \exp\!\left(-\left(\frac{Ma - 1}{w}\right)^{\!2}\right)
\]
With $C_{d,\infty} = 0.3$, $\Delta C_d = 0.5$, $w = 0.3$.\\
Maximum at $Ma=1$: $C_d = 0.3 + 0.5 = 0.8$ (almost 3$\times$ the supersonic value).
\end{theorybox}

% ──────────────────────────────────────────────────────────────
\section{Exercises}

\begin{exercisebox}{1 — Derive the Pressure Formula Step by Step}
Starting from:
\[
  \frac{dp}{dh} = -\frac{g\,p}{R_s\,(T_0 - Lh)}
\]
\textbf{a)} Separate variables:
\[
  \frac{dp}{p} = \underline{\hspace{6cm}}
\]

\textbf{b)} Integrate the right side using substitution $u = T_0 - Lh$ (so $du = -L\,dh$):
\[
  -\frac{g}{R_s}\int_{h=0}^{h} \frac{dh}{T_0 - Lh}
  = -\frac{g}{R_s}\int_{T_0}^{T(h)}\frac{du/(-L)}{u}
  = \frac{g}{R_s L}\int_{T(h)}^{T_0}\frac{du}{u}
  = \underline{\hspace{4cm}}
\]

\textbf{c)} Exponentiate both sides and apply $p(0) = p_0$ to arrive at the formula from Section 3.
\end{exercisebox}

\begin{solutionspace}[9cm]
\end{solutionspace}

\begin{exercisebox}{2 — Verify the Exponent Numerically}
Compute the exponent $\alpha = g/(R_s L)$ and $\alpha - 1$:
\[
  \alpha = \frac{9.80665}{287.058 \times 0.0065} = \underline{\hspace{3cm}}
\]
\[
  \alpha - 1 = \underline{\hspace{3cm}} \quad \text{(exponent for } \rho(h) \text{)}
\]

\textbf{Check:} Compute $\rho(5000)$ using both formulas and compare:
\begin{align*}
  T(5000) &= 288.15 - 0.0065 \times 5000 = \underline{\hspace{3cm}}\;\text{K} \\
  \rho(5000) &= 1.225 \times \left(\frac{T(5000)}{288.15}\right)^{4.256} = \underline{\hspace{3cm}}\;\text{kg/m}^3
\end{align*}
Reference value: $\rho(5000) \approx 0.7364\;\text{kg/m}^3$.
\end{exercisebox}

\begin{solutionspace}[6cm]
\end{solutionspace}

\begin{exercisebox}{3 — Compute the Atmosphere Table}
Fill in the complete table. Use the correct formula for each layer.

\begin{center}
\renewcommand{\arraystretch}{1.5}
\begin{tabular}{|c|c|c|c|c|}
\hline
$h$ (m) & $T$ (K) & $p$ (Pa) & $\rho$ (kg/m³) & $c$ (m/s) \\
\hline
0       & 288.15 & 101\,325 & 1.2250 & 340.3 \\
\hline
2\,000  & & & & \\
\hline
5\,000  & & & & \\
\hline
8\,000  & & & & \\
\hline
11\,000 & 216.65 & 22\,632 & 0.3639 & 295.1 \\
\hline
15\,000 & 216.65 & & & \\
\hline
20\,000 & 216.65 & 5\,475 & 0.0880 & \\
\hline
\end{tabular}
\end{center}
\end{exercisebox}

\begin{solutionspace}[6cm]
\end{solutionspace}

\begin{exercisebox}{4 — Mach Number at Altitude}
A rocket has speed $v = 600\;\text{m/s}$ throughout a vertical ascent.

\textbf{a)} Compute $Ma(h)$ at $h = 0$, $h = 5000\;\text{m}$, and $h = 11\,000\;\text{m}$:
\[
  Ma(0) = \frac{600}{\underline{\hspace{2cm}}} = \underline{\hspace{2cm}}, \quad
  Ma(5000) = \underline{\hspace{2cm}}, \quad
  Ma(11000) = \underline{\hspace{2cm}}
\]

\textbf{b)} The Mach number \textit{increases} even though the speed is constant. Explain why.

\textbf{c)} Compute $C_d$ at each Mach number above:
\[
  C_d(Ma) = 0.3 + 0.5\exp\!\left(-\left(\frac{Ma-1}{0.3}\right)^2\right)
\]
For $Ma = 1.033$: $C_d = 0.3 + 0.5\exp\left(-\left(\frac{0.033}{0.3}\right)^2\right) = \underline{\hspace{3cm}}$

\textbf{d)} Compute $F_\text{drag}$ at each altitude. Use $A = 10.73\;\text{m}^2$:
\[
  F_\text{drag}(0)\; = \tfrac{1}{2} \times \underline{\hspace{1cm}} \times 600^2 \times \underline{\hspace{1cm}} \times 10.73 = \underline{\hspace{2cm}}\;\text{N}
\]
\end{exercisebox}

\begin{solutionspace}[9cm]
\end{solutionspace}

\begin{exercisebox}{5 — The Barometric Formula: Approximation vs. ISA}
The simplified barometric formula (used in the 1D simulation) is:
\[
  \rho_\text{simple}(h) = \rho_0\,e^{-h/H}, \quad H = 8500\;\text{m}
\]

\textbf{a)} Derive this formula by solving $\frac{d\rho}{dh} = -\frac{\rho}{H}$ with $\rho(0) = \rho_0$ using separation of variables.

\textbf{b)} Compare $\rho_\text{simple}(h)$ vs.\ $\rho_\text{ISA}(h)$ at $h = 5000$, $10000$, $20000\;\text{m}$.
Compute the relative error $\dfrac{|\rho_\text{ISA} - \rho_\text{simple}|}{\rho_\text{ISA}}$ at each altitude.

\begin{center}
\begin{tabular}{|c|c|c|c|}
\hline
$h$ (m) & $\rho_\text{ISA}$ & $\rho_\text{simple}$ & Relative error \\
\hline
5\,000 & 0.7364 & & \\
\hline
10\,000 & 0.4135 & & \\
\hline
20\,000 & 0.0880 & & \\
\hline
\end{tabular}
\end{center}

\textbf{c)} Above what altitude does the simple model become unreliable (error $> 10\%$)?
\end{exercisebox}

\begin{solutionspace}[8cm]
\end{solutionspace}

\begin{exercisebox}{6 — Dynamic Pressure and Max-Q (Challenge)}
The \textbf{dynamic pressure} is $q = \frac{1}{2}\rho(h)\,v^2$.
During ascent, $\rho$ decreases while $v$ increases. The maximum of $q$ (called \textbf{Max-Q}) is a critical design point — structural loads on the rocket are highest here.

For a simplified vertical ascent with constant thrust $T = 1.4\;\text{MN}$, $m_0 = 100\,000\;\text{kg}$, $\dot{m} = 437.5\;\text{kg/s}$, assume $v \approx at$ (constant acceleration $a \approx 4.4\;\text{m/s}^2$).

\textbf{a)} Write $q(t) = \frac{1}{2}\rho(h(t))\,(at)^2$. With $h \approx \frac{1}{2}at^2$ and $\rho = \rho_0 e^{-h/H}$:
\[
  q(t) = \frac{1}{2}\rho_0\,e^{-at^2/(2H)} \cdot a^2 t^2
\]

\textbf{b)} Find the time $t^*$ of Max-Q by setting $\dfrac{dq}{dt} = 0$.
\textit{Hint: differentiate and factor. You will get $2 - \dfrac{a\,t^2}{H} = 0$.}
\[
  t^* = \sqrt{\frac{2H}{a}} = \underline{\hspace{3cm}}\;\text{s}
\]

\textbf{c)} Compute $h^*$ and $v^*$ at Max-Q. Does this match the Falcon 9 value of $\approx 80\;\text{s}$?
\end{exercisebox}

\begin{solutionspace}[9cm]
\end{solutionspace}

% ──────────────────────────────────────────────────────────────
\section{Resources}

\begin{hintbox}
\textbf{Videos:}
\begin{itemize}
  \item \textbf{Scott Manley} — ``Why does Max-Q matter?'':\\
        Search: ``Scott Manley Max-Q'' on YouTube
  \item \textbf{Everyday Astronaut} — ``What is Max-Q?'':\\
        \texttt{youtube.com/c/EverydayAstronaut} (excellent explanations with real launches)
  \item \textbf{Khan Academy} — ``Separable differential equations'':\\
        \texttt{khanacademy.org/math/ap-calculus-ab/ab-diff-equations-new/ab-7-6/v/separable-differential-equations-intro}
\end{itemize}
\textbf{Reference:}
\begin{itemize}
  \item ICAO Standard Atmosphere document: \texttt{engineeringtoolbox.com/international-standard-atmosphere-d\_985.html}
  \item HyperPhysics — Atmosphere: \texttt{hyperphysics.phy-astr.gsu.edu/hbase/Kinetic/barfor.html}
\end{itemize}
\end{hintbox}

\section{Connection to the Julia Template}

\begin{theorybox}{What you implement in \texttt{template.jl}}
\begin{enumerate}
  \item \texttt{isa\_temperature(h)}, \texttt{isa\_pressure(h)}, \texttt{isa\_density(h)}
  \item \texttt{speed\_of\_sound(h)}, \texttt{drag\_coefficient(mach)}, \texttt{aerodrag(h,v,A,Cd)}
  \item Plot all profiles $T(h)$, $\rho(h)$, $c(h)$ up to 20~km — overlay with simple exponential
  \item Dynamic pressure plot: $q(h)$ for a given velocity profile — find Max-Q visually
\end{enumerate}
\end{theorybox}

\end{document}
